\documentclass[11pt,a4paper]{article}
\usepackage[utf8]{inputenc}
\usepackage[T1]{fontenc}
\usepackage{lmodern}
\usepackage{textcomp}
\usepackage{amsmath,amssymb}
\usepackage{graphicx}
\usepackage{booktabs}
\usepackage{longtable}
\usepackage{array}
\usepackage{listings}
\usepackage{xcolor}
\usepackage[margin=1in]{geometry}
\usepackage{float}
\usepackage[hidelinks,breaklinks]{hyperref}
\usepackage{microtype}
\setlength{\parskip}{4pt}
\lstset{basicstyle=\ttfamily\scriptsize,breaklines=true,breakatwhitespace=false,
  postbreak=\mbox{\textcolor{gray}{$\hookrightarrow$}\space},frame=single,
  framerule=0.3pt,rulecolor=\color{gray!60},columns=fullflexible,keepspaces=true,
  showstringspaces=false,numbers=none,xleftmargin=4pt,xrightmargin=4pt,
  literate={\_}{{\_}}1}
\graphicspath{{figures/}}
\providecommand{\tightlist}{}
\title{Transcript: Viability of keV sterile-neutrino dark matter under BBN and X-ray constraints, 2024--2026}
\author{HEP-Genesis-Agent session record \\ \small Project: ramp-combo-sterile-bbn-xray}
\date{2026-10-21}
\begin{document}
\maketitle
\noindent\textit{\small The conversation verbatim, generated mechanically from the session record.}
\vspace{1em}



\appendix
\section{Conversation transcript}
\label{app:transcript}
The exchange verbatim (assistant text as shown to the user; tool calls summarized by name --- see the tool-call log for arguments and results).
\subsection*{User \small(2026-10-21 14:31:00)}
Survey recent arXiv work (2024-2026) on sterile-neutrino dark matter. Combine BBN constraints (use the sterile-nu tools for Yp/D/H numbers) and X-ray constraints: which (m\_s, sin\textasciicircum{}2 2theta) regions remain viable? Quantify with tool calls where possible.

\subsection*{Assistant \small(research mode, 2026-10-21 14:31:00)}
\textit{Plan:}\begin{enumerate}
  \item (closed) Use the sterile-nu tools to compute BBN (Yp, D/H, N\_eff) and relic-density constraints across the (m\_s, sin\textasciicircum{}2 2theta) plane. Combine them with the tool's X-ray line predictions to map which regions remain viable.
  \item (closed) Convert the 2024-2026 X-ray decay-rate limits (XRISM 2510.24560, with NuSTAR 1609.00667 as the high-mass reference) into sin\textasciicircum{}2 2theta(m\_s) upper limits, and cross-check them against the tool's X-ray predictions.
  \item (closed) Extract the quantitative 2024-2026 resonant-production (lepton-asymmetry) results: the largest L allowed by BBN and the CMB, the lower bound on sin\textasciicircum{}2 2theta, and the Lyman-alpha mass cutoff. These define the lower and low-mass edges of the viable band.
  \item (closed) Pin down the low-mass edge of the viable band from current structure-formation bounds (Lyman-alpha, Milky Way satellites, Tremaine-Gunn and phase-space limits), and map thermal-WDM mass bounds onto m\_s for resonant spectra.
  \item (blocked) Re-extract verbatim the numbers behind the contested findings f-001, f-002, f-004 and f-009. Resolve the m\_th fit inconsistency so the low-mass edge of the resonant window is set consistently, absorbing the structure-formation edge goal of d-04.
  \item (blocked) Quantify the upper (X-ray) and lower (asymmetry) edges of the viable band at 15, 20, 30, 40 and 60 keV, and assemble one viable-region table per asymmetry scenario.
  \item (closed) Robustness check of the BBN conclusions: test whether 'BBN does not bind keV sterile-neutrino DM' survives other error budgets, and estimate the Yp shift from lepton asymmetry that the tool omits.
\end{enumerate}
\textit{Tool calls:} \texttt{list\_sterile\_nu\_skills}, \texttt{list\_observables}, \texttt{list\_agent\_files}, \texttt{list\_input\_files}, \texttt{list\_datasets}, \texttt{list\_files}, \texttt{get\_dispatch}, \texttt{search\_arxiv\_abstracts}, \texttt{search\_arxiv\_abstracts}, \texttt{search\_arxiv\_abstracts}, \texttt{search\_arxiv}, \texttt{search\_arxiv\_abstracts}, \texttt{search\_arxiv\_abstracts}, \texttt{get\_arxiv\_metadata}, \texttt{get\_arxiv\_metadata}, \texttt{get\_arxiv\_metadata}, \texttt{search\_arxiv\_abstracts}, \texttt{search\_arxiv\_abstracts}, \texttt{get\_arxiv\_metadata}, \texttt{describe\_sterile\_nu\_tools}, \texttt{list\_sterile\_nu\_skills}, \texttt{load\_sterile\_nu\_skill}, \texttt{scan\_sterile\_nu\_parameters}, \texttt{download\_arxiv\_full\_text}, \texttt{download\_arxiv\_full\_text}, \texttt{download\_arxiv\_full\_text}, \texttt{predict\_sterile\_nu\_grid}, \texttt{download\_arxiv\_full\_text}, \texttt{predict\_sterile\_nu\_xray\_line}, \texttt{read\_text\_file}, \texttt{download\_arxiv\_full\_text}, \texttt{read\_arxiv\_chunk}, \texttt{read\_arxiv\_chunk}, \texttt{read\_arxiv\_chunk}, \texttt{read\_arxiv\_chunk}, \texttt{read\_arxiv\_chunk}, \texttt{read\_arxiv\_chunk}, \texttt{read\_arxiv\_chunk}, \texttt{read\_arxiv\_chunk}, \texttt{read\_arxiv\_chunk}, \texttt{read\_arxiv\_chunk}, \texttt{read\_arxiv\_chunk}, \texttt{predict\_sterile\_nu\_grid}, \texttt{predict\_sterile\_nu\_xray\_line}, \texttt{predict\_sterile\_nu\_xray\_line}, \texttt{predict\_sterile\_nu\_xray\_line}, \texttt{predict\_sterile\_nu\_xray\_line}, \texttt{predict\_sterile\_nu}, \texttt{post\_finding}, \texttt{post\_finding}, \texttt{post\_finding}, \texttt{post\_finding}, \texttt{post\_finding}, \texttt{post\_finding}, \texttt{post\_finding}, \texttt{post\_finding}, \texttt{plot\_sterile\_nu\_exclusion}, \texttt{predict\_sterile\_nu}, \texttt{cite}, \texttt{cite}, \texttt{cite}, \texttt{post\_finding}, \texttt{post\_finding}, \texttt{post\_finding}, \texttt{post\_finding}, \texttt{cite}, \texttt{cite}, \texttt{predict\_sterile\_nu\_xray\_line}, \texttt{post\_finding}, \texttt{post\_dead\_end}, \texttt{post\_question}, \texttt{post\_finding}, \texttt{get\_paper\_figures}, \texttt{read\_arxiv\_chunk}, \texttt{get\_paper\_figures}, \texttt{predict\_sterile\_nu\_grid}, \texttt{read\_arxiv\_chunk}, \texttt{read\_text\_file}, \texttt{read\_arxiv\_chunk}, \texttt{predict\_sterile\_nu\_xray\_line}, \texttt{predict\_sterile\_nu\_xray\_line}, \texttt{predict\_sterile\_nu\_xray\_line}, \texttt{predict\_sterile\_nu\_xray\_line}, \texttt{cite}, \texttt{read\_arxiv\_chunk}, \texttt{plot\_sterile\_nu\_bbn\_impact}, \texttt{predict\_sterile\_nu\_grid}, \texttt{post\_finding}, \texttt{post\_finding}, \texttt{post\_dead\_end}, \texttt{predict\_sterile\_nu\_xray\_line}, \texttt{predict\_sterile\_nu\_xray\_line}, \texttt{post\_finding}, \texttt{post\_question}, \texttt{post\_finding}, \texttt{post\_finding}, \texttt{post\_finding}, \texttt{post\_dead\_end}, \texttt{post\_finding}, \texttt{post\_finding}, \texttt{post\_finding}, \texttt{post\_finding}, \texttt{post\_question}\par
\textbf{Survey} --- 52 tools on 5 servers, 7 read-only probes, 9 papers found.

Warning: I could not use the sterile-neutrino BBN tools the task asks for. No such tools were available to me (only arXiv/web search and metadata lookup), so I have no computed Yp or D/H numbers. Every BBN statement below comes from paper abstracts, not from calculations. Literature state, 2024-2026: X-ray line searches keep tightening the limits and have found no new lines. The XRISM Collaboration stacked 3.75 Ms of Resolve data on ten galaxy clusters (2510.24560) and found no unidentified lines. Its 3-sigma limit on the decay rate of a 7.1 keV particle (the 3.55 keV line) is about 1.0e-27 s\textasciicircum{}-1. That is 3-4 times stronger than Hitomi's Perseus limit but still about 5 times weaker than the rate implied by the XMM-Newton stacked detection, so the 3.5 keV line is neither confirmed nor excluded. Older NuSTAR Galactic Center limits (1609.00667) remain the strongest for m\_s of 10-50 keV. Taken together, X-ray limits and structure-formation limits rule out non-resonant (Dodelson-Widrow) production as all of the dark matter, so the surviving parameter space needs resonant (Shi-Fuller) production driven by a lepton asymmetry. The main 2024-2026 development is on that BBN side. Full-neutrino-transport BBN/CMB bounds on lepton asymmetries (2405.06509) and quantum-kinetic studies of flavor equilibration (2502.14960) show that large flavor asymmetries are compatible with BBN and the CMB. This holds when they are present above roughly 20 MeV (Vogel et al.: L of order 1 is allowed) or when the total asymmetry nearly vanishes so that oscillations wash it out at MeV temperatures. This has reopened Shi-Fuller parameter space. Vogel et al. (2507.18752) find viable regions at m\_s \textgreater{}\textasciitilde{} 10 keV and sin\textasciicircum{}2 2theta \textless{}\textasciitilde{} 1e-14 for L \textgreater{}\textasciitilde{} 0.5, consistent with NuSTAR and INTEGRAL/SPI and with a Lyman-alpha preference for a thermal-equivalent WDM mass of about 4.1 keV. Akita, Hamaguchi and Ovchynnikov (2507.20659) widen the viable mixing-angle range by up to two orders of magnitude for m\_s up to 60 keV, and Gorbunov et al. (2502.17374) reopen lighter-mass regions via flavor-specific asymmetries and colder spectra. Qualitative picture of what survives: a band at roughly m\_s of 7-60 keV, with sin\textasciicircum{}2 2theta between the X-ray upper limit (falling roughly as m\_s\textasciicircum{}-5) and the lower bound set by the largest lepton asymmetry BBN and the CMB allow. The cut-off from structure formation sits at low mass, around m\_s of 5-10 keV depending on how the momentum spectrum is mapped. Note: the asymmetry-aware BBN windows rest on flavor-equilibration assumptions, and the CMB side is checked alongside BBN, so treat these windows as model-dependent.

Quantities worth checking: Radiative decay rate Gamma = 1.38e-29 s\textasciicircum{}-1 (sin\textasciicircum{}2 2theta / 1e-7)(m\_s / keV)\textasciicircum{}5, used to turn decay-rate limits into mixing-angle limits. My own estimate, not a tool result: the XRISM limit of 1e-27 s\textasciicircum{}-1 at 7.1 keV gives sin\textasciicircum{}2 2theta below about 4e-10, versus about 7e-11 for the XMM-Newton (Bulbul) line.; Yp and D/H as functions of the lepton asymmetry (xi\_e and the flavor split) and of Delta N\_eff; compare with Yp about 0.245 and D/H about 2.5e-5. This needs the BBN tools that were unavailable to me.; The largest total lepton asymmetry BBN and the CMB allow; it sets the lowest mixing angle at which resonant production can still supply all the dark matter.; Thermal-equivalent WDM mass for each (m\_s, sin\textasciicircum{}2 2theta, L) point versus Lyman-alpha bounds of about 4-6 keV.; Required L(m\_s, sin\textasciicircum{}2 2theta) to give Omega\_s h\textasciicircum{}2 = 0.12, checked against the power-law fits in 2507.18752..

\textbf{Plan} --- 4 lines of investigation:

\begin{enumerate}
  \item Use the sterile-nu tools to compute BBN (Yp, D/H, N\_eff) and relic-density constraints across the (m\_s, sin\textasciicircum{}2 2theta) plane. Combine them with the tool's X-ray line predictions to map which regions remain viable. --- The task asks for tool-computed Yp and D/H numbers. The survey could not obtain any, so this direction answers the question most directly.
  \item Convert the 2024-2026 X-ray decay-rate limits (XRISM 2510.24560, with NuSTAR 1609.00667 as the high-mass reference) into sin\textasciicircum{}2 2theta(m\_s) upper limits, and cross-check them against the tool's X-ray predictions. --- X-ray limits set the upper edge of the viable band. Reading the limits verbatim from the papers lets us validate both the survey's estimate (about 4e-10 at 7.1 keV) and the tool's numbers.
  \item Extract the quantitative 2024-2026 resonant-production (lepton-asymmetry) results: the largest L allowed by BBN and the CMB, the lower bound on sin\textasciicircum{}2 2theta, and the Lyman-alpha mass cutoff. These define the lower and low-mass edges of the viable band. --- If the tools model only Dodelson-Widrow production, the surviving Shi-Fuller window has to come from the literature. The papers' BBN numbers can also be cross-checked against the tool's Yp and D/H.
  \item Pin down the low-mass edge of the viable band from current structure-formation bounds (Lyman-alpha, Milky Way satellites, Tremaine-Gunn and phase-space limits), and map thermal-WDM mass bounds onto m\_s for resonant spectra. --- The brief's low-mass cutoff (5-10 keV, Lyman-alpha 'preference' at 4.1 keV) is unsourced, and d-03 has no budget left to cover it. Recent bounds (e.g. Villasenor 2023, Irsic 2024, DES satellites) are much stronger than 4 keV and could close the 7 keV region. (added by the critic)
Critic's view: revise --- Step 1 must check the tool's validity range and production mechanism before the scan is configured. If the tool models only Dodelson-Widrow production, a 1e-14 to 1e-6 grid is mostly wasted, and the resonant window cannot be called tool-computed. Make the scan bounds depend on what describe\_sterile\_nu\_tools reports.; A keV-mass sterile neutrino that is non-relativistic at BBN and has tiny mixing gives Delta N\_eff of about 0, so Yp and D/H will sit at their Standard Model values. Do not present that as a BBN constraint on the plane. State whether the tool's BBN output responds to m\_s and sin\textasciicircum{}2 2theta at all, and give the size of the deviation alongside its observational uncertainty (Yp about 0.245 +/- 0.003; D/H about 2.53 +/- 0.03 e-5). If the tool accepts a lepton asymmetry, run a small xi\_e scan instead, since that is where BBN actually bites.; No general plotting tool exists, so a combined literature-plus-tool overlay cannot be produced. Report only the plot\_sterile\_nu\_exclusion and plot\_sterile\_nu\_bbn\_impact paths, and give the combined viable band as a table. check\_sterile\_nu\_sbl\_anomaly is off-task and should be dropped if the budget is tight..

\end{enumerate}
\textbf{Wave 1} --- now exploring, in parallel: Use the sterile-nu tools to compute BBN (Yp, D/H, N\_eff) and relic-density constraints across the (m\_s, sin\textasciicircum{}2 2theta) plane. Combine them with the tool's X-ray line predictions to map which regions remain viable.; Convert the 2024-2026 X-ray decay-rate limits (XRISM 2510.24560, with NuSTAR 1609.00667 as the high-mass reference) into sin\textasciicircum{}2 2theta(m\_s) upper limits, and cross-check them against the tool's X-ray predictions.; Extract the quantitative 2024-2026 resonant-production (lepton-asymmetry) results: the largest L allowed by BBN and the CMB, the lower bound on sin\textasciicircum{}2 2theta, and the Lyman-alpha mass cutoff. These define the lower and low-mass edges of the viable band..

\textbf{Finding:} Vogel et al. (2507.18752): Shi-Fuller sterile-nu DM is viable for initial L \textgreater{}\textasciitilde{} 0.5 (L given at T=10 GeV), m\_s \textgreater{}\textasciitilde{} 35 keV and sin\textasciicircum{}2 2theta \textless{}\textasciitilde{} 1e-14 once X-ray limits (NuSTAR+INTEGRAL/SPI, 95\% CL) and the structure-formation bound m\_th \textgreater{} 8.5 keV (Zelko+, 95\% CL) are applied; the SF window is excluded at \textgreater{}\textgreater{}95\% CL for L \textless{} 0.1; their Fig. 7 covers m\_s 15-60 keV, sin\textasciicircum{}2 2theta 1e-15 to 1e-13.

\textbf{Finding:} 2507.18752 provides fitting functions valid for 15 keV \textless{}= m\_s \textless{}= 60 keV and 1e-15 \textless{}= sin\textasciicircum{}2 2theta \textless{}= 3e-13: m\_th = 1.329 m\_s\textasciicircum{}1.044 (sin\textasciicircum{}2 2theta)\textasciicircum{}0.055 (keV; average error 1.3\%) and the L needed for all DM, L = 1.865e-12 m\_s\textasciicircum{}-1.813 (sin\textasciicircum{}2 2theta)\textasciicircum{}-1.042 (average error 8.3\%); also the mapping m\_th,new = 0.83 m\_th,old + 0.36 keV. Lyman-alpha (Villasenor+) prefers m\_th = 4.1 (+33.3/-1.5) keV and the whole Fig. 7 region is within that 1-sigma range; galaxy counts give m\_th \textgreater{} 8.4 keV, and counts plus lensing give \textgreater{} 8.5 keV (95\% CL). My own evaluation of the fit, not a tool result: at the paper's benchmark m\_s = 37 keV, sin\textasciicircum{}2 2theta = 1.5e-14, it gives L \textasciitilde{} 0.68 and m\_th \textasciitilde{} 8.6 keV.

\textbf{Finding:} Akita, Hamaguchi, Ovchynnikov (2507.20659): with large flavor asymmetries and zero total asymmetry (washed out by oscillations at T\_osc \textasciitilde{} 15 MeV), the viable sin\textasciicircum{}2 2theta range grows by up to two orders of magnitude for m\_s up to 60 keV. Individual asymmetries up to \textbar{}L\_alpha\textbar{} \textasciitilde{} 0.1 are allowed by BBN/CMB (benchmarks L\_e = -L\_mu = 0.1 and L\_mu = -L\_tau = 0.06). The largest flavor-universal asymmetry allowed by BBN+CMB is L\_e = L\_mu = L\_tau = 1e-3. Their Fig. 1 plane spans sin\textasciicircum{}2 2theta 1e-16 to 1e-8; the lower edge corresponds to sum L\_alpha = 0. Lyman-alpha is shown only as an illustrative guide. Public code: sterile-dm-lfa on GitHub.

\textbf{Finding:} Model dependence: Froustey and Pitrou (2405.06509), using full 3-flavor neutrino transport plus BBN, find that 'the space of allowed initial asymmetries is generically unbound' with current abundance and CMB data, because flavor equilibration is not complete in general and the CMB bound is at least quadratic in the asymmetries. So there is no single L\_max. The L \textgreater{}\textasciitilde{} 0.5-1 (2507.18752) and \textbar{}L\_alpha\textbar{} \textasciitilde{} 0.1 with sum = 0 (2507.20659) windows depend on the initial flavor configuration at T \textgreater{}\textasciitilde{} 20 MeV. Only for flavor-universal asymmetry is the bound tight (about 1e-3).

\textbf{Finding:} Sterile-nu tools: BBN uses a PRyMordial fit (Yp, D/H, Li7) driven only by Delta N\_eff, with exclusion at Delta chi\textasciicircum{}2 \textgreater{} 3.84 against Yp = 0.2449 +/- 0.004 (Aver+ 2021), D/H = 2.547 +/- 0.025 e-5 (Cooke+ 2018) and Planck N\_eff = 2.99 +/- 0.17. The model is valid for m\_s below about 100 keV. Production is DW; there is an SF option with a lepton\_asymmetry parameter, but its Delta N\_eff is a heuristic (0.1 x DW) and is not a resonant-production calculation. The tools do not model the effect of the asymmetry itself on BBN (xi\_e).

\textbf{Finding:} DW scan (m\_s 1-100 keV, sin\textasciicircum{}2 2theta 1e-14 to 1e-6, 60x60): 546 of 3600 points are BBN-excluded and 1740 overclose. The BBN 95\% upper limit on sin\textasciicircum{}2 2theta is 3.9e-7 at 1 keV, 8.2e-8 at 7 keV, 4.4e-8 at 13 keV, 2.4e-8 at 26 keV, 1.7e-8 at 50 keV and 9.3e-9 at 100 keV. This is two or more orders of magnitude weaker than the overclosure boundary, so BBN does not constrain keV sterile-neutrino DM.

\textbf{Finding:} Tool BBN numbers (DW, 84 points): the SM baseline is Yp = 0.246815, D/H = 2.4433e-5 and N\_eff = 3.044. This is +0.48 sigma in Yp and -4.1 sigma in D/H from the observed values (chi\textasciicircum{}2 = 2.77). On the DW relic curve (Omega h\textasciicircum{}2 = 0.12: sin\textasciicircum{}2 2theta = 2.545e-10 at 7.1 keV, 1.374e-10 at 10 keV, 3.946e-11 at 20 keV, 7.583e-12 at 50 keV), Delta N\_eff = 1.6e-3, 1.1e-3, 5.6e-4 and 2.2e-4, and Delta chi\textasciicircum{}2\_BBN = -0.020, -0.014, -0.007 and -0.003. At 7.1 keV, sin\textasciicircum{}2 2theta = 3e-10 gives Yp = 0.246840 and D/H = 2.44395e-5. BBN starts to exclude only at sin\textasciicircum{}2 2theta around 1e-7 (7.1 keV: Delta N\_eff = 0.62, Yp = 0.2549, D/H = 2.649e-5, Delta chi\textasciicircum{}2 = 5.9) or around 3e-8 (50 keV: Delta N\_eff = 0.88, Delta chi\textasciicircum{}2 = 16.8).

\textbf{Finding:} X-ray decay on the DW relic curve: Gamma = 6.25e-28 s\textasciicircum{}-1 (7.1 keV, E = 3.55 keV), 1.87e-27 (10 keV), 1.72e-26 (20 keV) and 3.23e-25 (50 keV). These agree within about 1.5\% with the analytic 1.38e-29 (sin\textasciicircum{}2 2theta/1e-7)(m\_s/keV)\textasciicircum{}5 formula (6.34e-28 at 7.1 keV, 3.27e-25 at 50 keV); the tool uses a 1.36e-29 prefactor. Tabulated limits exclude the DW-DM points at 7.1 and 10 keV (NuSTAR M31 Roach+2020 and XMM blank-sky Dessert+2020) and at 20 keV (NuSTAR M31). The 50 keV point is not excluded only because the tool has no tabulated limit within 30\% in mass. Whole-halo Milky Way-center fluxes are 11.9, 25.3, 116 and 873 ph/s/cm\textasciicircum{}2 (no field-of-view cut). DW production as all of the DM is therefore X-ray excluded for 7-20 keV.

\textbf{Finding:} XRISM 2510.24560 (stacked spectra of 10 clusters, 3.75 Ms) gives a 3-sigma upper limit of 0.95e-6 ph/cm\textasciicircum{}2/s on line flux at the 3.55 keV line, which converts to a decay-rate limit of Gamma \textless{} 0.97e-27 s\textasciicircum{}-1 for m\_s = 7.1 keV. The paper's mixing-angle fit, sin\textasciicircum{}2 2theta/1e-10 \textless{} 4.1 (E/3.55)\textasciicircum{}-4 f1(E), with f1 = 7.8774 - 3.5017E + 0.59087E\textasciicircum{}2 - 0.042245E\textasciicircum{}3 + 0.0010943E\textasciicircum{}4 (E in keV, valid 2.5-15 keV), evaluates by hand to the following upper limits on sin\textasciicircum{}2 2theta: about 1.3e-9 at m\_s = 6 keV, about 4-5e-10 at 7.1 keV (3.9e-10 from Gamma; 4.8e-10 from the polynomial, f1 = 1.18), about 2.3e-10 at 8 keV, about 5.7e-11 at 10 keV, and about 1.2e-11 at 14 keV.

\textbf{Finding:} The tool's X-ray decay rate at m\_s = 7.1 keV with sin\textasciicircum{}2 2theta = 4e-10 is 9.82e-28 s\textasciicircum{}-1 (lifetime 1.02e27 s). This agrees to within 2\% with the XRISM limit of 0.97e-27 s\textasciicircum{}-1, so the tool's rate normalization matches the standard formula Gamma = 1.38e-32 s\textasciicircum{}-1 (sin\textasciicircum{}2 2theta/1e-10)(m/keV)\textasciicircum{}5 quoted in NuSTAR 1609.00667 Eq. 2. The tool also flags this point as excluded by NuSTAR M31 (Roach+ 2020) and XMM blank-sky (Dessert+ 2020), both stronger than XRISM. The tool's fluxes are whole-halo values (f\_DM = 1, no field-of-view cut), so they cannot be compared directly with field-of-view-limited published flux limits; only the decay rate can be compared.

\textbf{Finding:} The NuSTAR Galactic Center paper (1609.00667) gives its sterile-neutrino line flux as F of about 2.6e-6 cm\textasciicircum{}-2 s\textasciicircum{}-1 (m/20 keV)\textasciicircum{}4 (sin\textasciicircum{}2 2theta/1e-14)(DeltaOmega/4 deg\textasciicircum{}2)(J/40 GeV cm\textasciicircum{}-3 kpc sr\textasciicircum{}-1), using the coreNFW default with J\_tot = 29 and DeltaOmega\_tot of about 3.8 deg\textasciicircum{}2. It states the limit improves on earlier ones by about an order of magnitude near m\_s = 20 keV and implies m\_s \textless{}\textasciitilde{} 16 keV in the nuMSM. The paper gives its sin\textasciicircum{}2 2theta limits at 10, 20 and 40 keV only as curves in Figs. 7-8, not as numbers in the text, so the full-text extraction yields no verbatim values at those masses.

\textbf{Finding:} I cross-checked the tool against the 2507.18752 benchmark point, with SF production at m\_s = 37 keV, sin\textasciicircum{}2 2theta = 1.5e-14 and L = 0.1 (the tool maximum). The tool gives Delta N\_eff = 3.5e-8, Yp = 0.24682, D/H = 2.443e-5, Delta chi\textasciicircum{}2(BBN) = 0.00 (allowed), Gamma = 1.42e-28 s\textasciicircum{}-1 (E = 18.5 keV), and Omega\_s h\textasciicircum{}2 = 25.1 (overcloses). The BBN output equals the SM value: the asymmetry does not enter the BBN abundances (no xi\_e shift of n/p). Its SF Delta N\_eff is a heuristic (0.1 x DW). The tool's SF relic density is about 200 times the observed DM, whereas the paper's fit puts this point at the correct DM density for L \textasciitilde{} 0.7. So the tool cannot test the published resonant windows: neither the asymmetry-dependent BBN numbers nor the SF relic density.

\textbf{Finding:} DW exclusion-plane plot for m\_s 1-100 keV, sin\textasciicircum{}2 2theta 1e-14 to 1e-6 (Figure 10, exclusion\_scan\_DW\_9bfc3d.png). Viable DW region: the Omega h\textasciicircum{}2 = 0.12 line runs from sin\textasciicircum{}2 2theta = 2.5e-10 at 7.1 keV to 7.6e-12 at 50 keV, with overclosure above it. BBN is irrelevant (it bites only above about 1e-8 to 1e-7). X-ray limits exclude the DW relic line for 7-20 keV. Below the line, sub-dominant DW sterile neutrinos are allowed, down to the X-ray limits. Full-DM viability needs resonant production, which the tool does not model reliably.

\textbf{Finding:} Table of upper limits on sin\textasciicircum{}2 2theta from X-rays. From XRISM 2510.24560 (3-sigma limit, my arithmetic on the paper's polynomial): m\_s = 6 keV, about 1.3e-9; 7.1 keV, about 4-5e-10 (3.9e-10 from Gamma \textless{} 0.97e-27 s\textasciicircum{}-1); 8 keV, about 2.3e-10; 10 keV, about 5.7e-11; 14 keV, about 1.2e-11. NuSTAR GC 1609.00667: improves on earlier limits by about 10x near 20 keV, but the numbers exist only in Figs. 7-8, so none are quoted. Tool check at m\_s = 20 keV, sin\textasciicircum{}2 2theta = 1e-12: Gamma = 4.36e-28 s\textasciicircum{}-1, versus 4.42e-28 s\textasciicircum{}-1 from 1609.00667 Eq. 2 (agreement 1.4\%). The tool flags this point as excluded by NuSTAR M31 (Roach+ 2020), so the tool's tabulated limit at 20 keV lies below 1e-12.

\textbf{Completed --- Use the sterile-nu tools to compute BBN (Yp, D/H, N\_eff) and relic-density constraints across the (m\_s, sin\textasciicircum{}2 2theta) plane. Combine them with the tool's X-ray line predictions to map which regions remain viable..} The sterile-nu tools give BBN numbers only through Delta N\_eff (a PRyMordial fit for Yp, D/H and Li7; exclusion at Delta chi\textasciicircum{}2 \textgreater{} 3.84 against Yp 0.2449 +/- 0.004, D/H 2.547 +/- 0.025 e-5 and Planck N\_eff 2.99 +/- 0.17; valid below about 100 keV). Production is DW only in any reliable sense. The SF mode accepts a lepton\_asymmetry but is a heuristic (0.1 x DW for Delta N\_eff), and it returned an unphysical Omega h\textasciicircum{}2 = 432 at 7.1 keV, sin\textasciicircum{}2 2theta = 7e-11, L = 1e-3. The lepton asymmetry's own effect on BBN is not modeled (f-005, dead end x-01; answers q-01). DW scan over m\_s 1-100 keV and sin\textasciicircum{}2 2theta 1e-14 to 1e-6 (60x60, scan\_DW\_9bfc3d.npz/.csv): 546 of 3600 points are BBN-excluded and 1740 overclose. The BBN 95\% upper limit on sin\textasciicircum{}2 2theta is 8.2e-8 at 7 keV and 1.7e-8 at 50 keV, two or more orders of magnitude weaker than overclosure, so BBN does not constrain keV DM (f-006; answers q-04). The tool tabulated Yp, D/H and Delta N\_eff at 84 points (sterile\_nu\_grid\_DW\_097459.csv, sterile\_nu\_grid\_DW\_1b6ff9.csv). The SM baseline is Yp = 0.24682 (+0.48 sigma) and D/H = 2.443e-5 (-4.1 sigma, a baseline offset; question q-05). On the DW relic curve (sin\textasciicircum{}2 2theta = 2.545e-10, 1.374e-10, 3.946e-11 and 7.583e-12 at 7.1, 10, 20 and 50 keV), Delta N\_eff is 1.6e-3 to 2.2e-4 and \textbar{}Delta chi\textasciicircum{}2\textbar{} is 0.02 or less. BBN starts to exclude only at sin\textasciicircum{}2 2theta of about 1e-7 at 7.1 keV (Delta N\_eff = 0.62, Yp = 0.2549, D/H = 2.649e-5) (f-007). X-ray decay rates on the relic curve are 6.25e-28, 1.87e-27, 1.72e-26 and 3.23e-25 s\textasciicircum{}-1, within about 1.5\% of the analytic 1.38e-29 formula. The DW relic line is excluded at 7.1, 10 and 20 keV by NuSTAR M31 and XMM blank sky. At 50 keV the tool has no tabulated limit within 30\% in mass (f-008). Fluxes are whole-halo, with no field-of-view cut. Exclusion-plane plot: Figure 10, /Users/nramachandra/hep-genesis-testing/projects/ramp-combo-sterile-bbn-xray/runs/261021\_143100/exclusion\_scan\_DW\_9bfc3d.png (f-013). Viable region, DW: sub-dominant DM below the Omega h\textasciicircum{}2 = 0.12 line and below the X-ray limits. DW as all of the DM is excluded for 7-20 keV, so full-DM viability needs resonant production, which must come from the literature (d-03). Not done: plot\_sterile\_nu\_bbn\_impact and check\_sterile\_nu\_sbl\_anomaly were not run because the budget ran out.

\textbf{Completed --- Extract the quantitative 2024-2026 resonant-production (lepton-asymmetry) results: the largest L allowed by BBN and the CMB, the lower bound on sin\textasciicircum{}2 2theta, and the Lyman-alpha mass cutoff. These define the lower and low-mass edges of the viable band..} Main result: the published resonant-production window survives only at large lepton asymmetry. It needs L \textgreater{}\textasciitilde{} 0.5 (set at T = 10 GeV), m\_s \textgreater{}\textasciitilde{} 35 keV and sin\textasciicircum{}2 2theta \textless{}\textasciitilde{} 1e-14 once X-ray limits and the structure bound m\_th \textgreater{} 8.5 keV (95\% CL) are applied. For L \textless{} 0.1 the window is excluded at much more than 95\% CL. The sterile-neutrino tool cannot test this window.

There is no single largest L allowed by BBN and the CMB. Froustey and Pitrou (2405.06509) find the allowed initial asymmetries 'generically unbound' with current abundance and CMB data, because flavors do not fully equilibrate. The bound depends on the flavor pattern:

\begin{itemize}
  \item Same asymmetry in every flavor: L\_e = L\_mu = L\_tau \textless{}\textasciitilde{} 1e-3 (cited in 2507.20659).
  \item Opposite flavor asymmetries summing to zero: \textbar{}L\_alpha\textbar{} up to about 0.1 (2507.20659). These wash out by oscillations below about 15 MeV.
  \item Large asymmetry above 20 MeV: L of order 1 (2507.18752). Their BBN check uses L\_mu about 0.68 at 20 MeV, and L drops by a factor of 5-10 by 10 MeV.
\end{itemize}
Vogel et al. (2507.18752):

\begin{itemize}
  \item Fig. 7 covers m\_s 15-60 keV and sin\textasciicircum{}2 2theta 1e-15 to 3e-13.
  \item Fitting functions (valid for 15-60 keV and 1e-15 to 3e-13): L = 1.865e-12 m\_s\textasciicircum{}-1.813 (sin\textasciicircum{}2 2theta)\textasciicircum{}-1.042, with 8.3\% average error; and m\_th = 1.329 m\_s\textasciicircum{}1.044 (sin\textasciicircum{}2 2theta)\textasciicircum{}0.055 keV, with 1.3\% average error. Their Eq. 7 converts older thermal-WDM masses: m\_th,new = 0.83 m\_th,old + 0.36 keV.
  \item Lyman-alpha prefers a thermal-WDM mass of 4.1 (+33.3/-1.5) keV, and the whole Fig. 7 region lies within that 1-sigma range. Galaxy counts give m\_th \textgreater{} 8.4 keV; counts plus lensing give \textgreater{} 8.5 keV (95\% CL). This lensing-plus-counts bound sets the m\_s \textgreater{}\textasciitilde{} 35 keV low-mass edge.
  \item My own hand evaluation of the fits, not a tool or paper number: at the paper's benchmark (37 keV, sin\textasciicircum{}2 2theta = 1.5e-14) they give L about 0.68 and m\_th about 8.6 keV.
\end{itemize}
Lower bound on sin\textasciicircum{}2 2theta at two or more masses: I could not get these as numbers. The lower edge of the viable region in 2507.20659 Fig. 1 corresponds to zero total asymmetry. They report mixing angles extending by up to two orders of magnitude for m\_s up to 60 keV, on a plot spanning 1e-16 to 1e-8. Their code is public (sterile-dm-lfa on GitHub), and the Vogel et al. code is at github.com/garciaeh77/GZA. In 2507.18752 the lower edge in the 15-60 keV range comes from the L fit above.

Tool cross-check (resonant production, m\_s = 37 keV, sin\textasciicircum{}2 2theta = 1.5e-14, L = 0.1, the tool's maximum): Delta N\_eff = 3.5e-8, Yp = 0.24682, D/H = 2.443e-5, BBN Delta chi\textasciicircum{}2 = 0 (allowed), decay rate 1.42e-28 s\textasciicircum{}-1 at 18.5 keV, and Omega\_s h\textasciicircum{}2 = 25.1 (overcloses). The asymmetry has no effect on the BBN numbers, and the resonant Delta N\_eff is a placeholder (0.1 times the Dodelson-Widrow value). The relic density is about 200 times the observed dark matter, while the paper puts this point at the observed density for L about 0.7. So the tool cannot test the resonant windows.

Not done: I did not read 2502.17374 (reopened low-mass region) or extract the parameter-space figures with get\_paper\_figures; I ran out of budget.

Files: findings f-001 to f-004 and f-012; full texts 2405.06509.txt, 2507.18752.txt and 2507.20659.txt, plus references.bib (with cited quotes for all three papers), all in /Users/nramachandra/hep-genesis-testing/projects/ramp-combo-sterile-bbn-xray/runs/261021\_143100/.

\textbf{Completed --- Convert the 2024-2026 X-ray decay-rate limits (XRISM 2510.24560, with NuSTAR 1609.00667 as the high-mass reference) into sin\textasciicircum{}2 2theta(m\_s) upper limits, and cross-check them against the tool's X-ray predictions..} I turned the XRISM 2510.24560 limits into upper limits on sin\textasciicircum{}2 2theta and checked the tool against them. One gap remains: I have no NuSTAR mixing-angle values at 10, 20 or 40 keV. 1609.00667 gives those limits only as curves in its Figs. 7-8, and I did not run get\_paper\_figures or digitize the plots before the budget ran out.

XRISM, cited verbatim: the 3-sigma limit on the decay rate at m\_s = 7.1 keV is Gamma \textless{} 0.97e-27 s\textasciicircum{}-1. The paper also gives a mixing-angle fit, sin\textasciicircum{}2 2theta/1e-10 \textless{} 4.1(E/3.55)\textasciicircum{}-4 f1(E), valid for E = 2.5-15 keV.

Upper limits on sin\textasciicircum{}2 2theta from that fit (my own arithmetic on the paper's formula, not a tool result):

\begin{itemize}
  \item 6 keV: about 1.3e-9
  \item 7.1 keV: about 4-5e-10 (3.9e-10 from the Gamma limit, 4.8e-10 from the fit). This confirms the survey's estimate of about 4e-10.
  \item 8 keV: about 2.3e-10
  \item 10 keV: about 5.7e-11
  \item 14 keV: about 1.2e-11
\end{itemize}
The 3.5 keV line is not excluded: its implied rate of about 2e-28 s\textasciicircum{}-1 (sin\textasciicircum{}2 2theta about 7e-11) sits a factor of 5 below the XRISM limit.

NuSTAR, cited verbatim: the paper quotes a decay-rate formula (Eq. 2), Gamma = 1.38e-32 s\textasciicircum{}-1 (sin\textasciicircum{}2 2theta/1e-10)(m/keV)\textasciicircum{}5. It states that near 20 keV it improves earlier limits by about an order of magnitude and implies m\_s \textless{}\textasciitilde{} 16 keV in the nuMSM.

Tool cross-check:

\begin{itemize}
  \item At 7.1 keV with sin\textasciicircum{}2 2theta = 4e-10, the tool gives Gamma = 9.82e-28 s\textasciicircum{}-1, within 2\% of the XRISM limit. The tool also marks this point as excluded by NuSTAR M31 (Roach+ 2020) and XMM blank-sky (Dessert+ 2020).
  \item At 20 keV with sin\textasciicircum{}2 2theta = 1e-12, the tool gives Gamma = 4.36e-28 s\textasciicircum{}-1, against 4.42e-28 s\textasciicircum{}-1 from Eq. 2. The tool marks this point as excluded by NuSTAR M31, so its limit at 20 keV is below 1e-12.
\end{itemize}
Normalization (answers q-02): the tool's fluxes are for the whole halo with no field-of-view cut, so only decay rates can be compared with the papers. The papers use different geometry factors: a projected-mass w factor for XRISM and J\_tot = 29 GeV cm\textasciicircum{}-3 kpc sr\textasciicircum{}-1 over DeltaOmega of about 3.8 deg\textasciicircum{}2 for NuSTAR.

Findings: f-009, f-010, f-011, f-014. References: /Users/nramachandra/hep-genesis-testing/projects/ramp-combo-sterile-bbn-xray/runs/261021\_143100/references.bib. Full texts: 2510.24560.txt and 1609.00667.txt in the same directory.

\textbf{Critic, wave 1} --- continue: 9 findings stand up to the evidence, 5 questioned ('Vogel et al. (2507.18752): Shi-Fuller sterile-nu DM is viable for initial L \textgreater{}\textasciitilde{} 0.5 (L give': Neither excerpt contains the claimed values. tc-35 is the title page and tc-40 is Fig. 4 text about 20 MeV asymmetry configurations. Missing: L \textgreater{}\textasciitilde{} 0.5 at 10 GeV, m\_s \textgreater{}\textasciitilde{} 35 keV, sin\textasciicircum{}2 2theta \textless{}\textasciitilde{} 1e-14, m\_th \textgreater{} 8.5 keV (Zelko+) and the \textgreater{}\textgreater{}95\% CL exclusion for L \textless{} 0.1. The Fig. 7 range (1e-15 to 1e-13) conflicts with the fit range in f-002 (up to 3e-13). The board brief says m\_s \textgreater{}\textasciitilde{} 10 keV for the same paper, which conflicts with 35 keV.; '2507.18752 provides fitting functions valid for 15 keV \textless{}= m\_s \textless{}= 60 keV and 1e-15 \textless{}= sin\textasciicircum{}2': The excerpts (the same as f-001) show none of the fit coefficients, the m\_th mapping, or the Villasenor 4.1 (+33.3/-1.5) keV value. Re-evaluating the stated fits at 37 keV and sin\textasciicircum{}2 2theta = 1.5e-14 gives L of about 0.68, which agrees. It gives m\_th = 1.329 x 37\textasciicircum{}1.044 x (1.5e-14)\textasciicircum{}0.055, about 10.0 keV, not the claimed 8.6 keV. So either the transcribed coefficients or the hand arithmetic are wrong, and this also moves the m\_s \textgreater{}\textasciitilde{} 35 keV edge.; 'Model dependence: Froustey and Pitrou (2405.06509), using full 3-flavor neutrino transport': Both excerpts are title and author blocks. The quoted phrase 'generically unbound' and the reasoning that the CMB bound is at least quadratic do not appear in the evidence.; 'XRISM 2510.24560 (stacked spectra of 10 clusters, 3.75 Ms) gives a 3-sigma upper limit of ': Excerpts tc-37 (affiliations) and tc-41 (NXB modelling) do not contain the 0.95e-6 ph/cm\textasciicircum{}2/s flux limit, the Gamma \textless{} 0.97e-27 s\textasciicircum{}-1 limit, or the f1(E) polynomial. Recomputing from the stated polynomial reproduces the hand values (f1(3.55) of about 1.18 gives 4.8e-10, and 5.7e-11 at 10 keV), so the arithmetic is consistent. The source values still need a verbatim excerpt.).

Issues: The baseline numbers are internally inconsistent. Offsets of +0.48 sigma in Yp and -4.15 sigma in D/H (sigma = 0.025e-5) give chi\textasciicircum{}2 of about 17.4, not the tool's 2.77. The tool's likelihood must include an extra D/H theory error, so the '-4.1 sigma' statement misdescribes it. The negative Delta chi\textasciicircum{}2 values and the Delta N\_eff exclusion threshold depend on this error budget.; The 1.5\% gap from the analytic formula comes mostly from display rounding (2.5e-10 printed, 2.545e-10 used). The missing 50 keV limit is a gap in the tool, not a physical allowance. Published NuSTAR GC, INTEGRAL and Lyman-alpha limits exclude DW-DM at 50 keV, and the write-up should say so.; The DW viable region relies on the tool's tabulated X-ray limits (Roach+ 2020, Dessert+ 2020). Their numerical sin\textasciicircum{}2 2theta curves were never extracted, so 'below the X-ray limits' is not quantified. The XMM blank-sky limit itself is contested in the literature (Boyarsky et al. vs Dessert et al.), and the robustness of the conclusion to that choice was not checked..

\textbf{Next steps after wave 1} --- This is the last wave (13 calls, about 24 minutes, at most 3 directions), so it goes to assembling a defensible final answer rather than new exploration. I close d-04 because its goal, the structure-formation low-mass edge, is now covered by the first direction: that edge turns on the m\_th \textgreater{} 8.5 keV bound and the disputed m\_th fit (10.0 vs 8.6 keV at the benchmark). A broad new Lyman-alpha or satellite survey would not fit the budget. The first direction answers the critic's high-priority request: verbatim support for the contested findings f-001, f-002, f-004, f-009 and f-011, mostly read from full texts already downloaded, which keeps the call count low. The second is the core deliverable and the cross-check against the literature: it digitizes the published X-ray and asymmetry edges at 15-60 keV and checks them against a single tool grid call, giving a viable-interval table per asymmetry scenario. The third, lowest priority, is a robustness check of the BBN conclusion under other likelihood error budgets, plus a quantitative estimate of the lepton-asymmetry effect on Yp, which the tool does not model. The task does not ask for facility runs, so none is scheduled. Planned use is about 10 tool calls.

New lines of investigation: Re-extract verbatim the numbers behind the contested findings f-001, f-002, f-004 and f-009. Resolve the m\_th fit inconsistency so the low-mass edge of the resonant window is set consistently, absorbing the structure-formation edge goal of d-04.; Quantify the upper (X-ray) and lower (asymmetry) edges of the viable band at 15, 20, 30, 40 and 60 keV, and assemble one viable-region table per asymmetry scenario.; Robustness check of the BBN conclusions: test whether 'BBN does not bind keV sterile-neutrino DM' survives other error budgets, and estimate the Yp shift from lepton asymmetry that the tool omits..

\textbf{Wave 2} --- now exploring, in parallel: Re-extract verbatim the numbers behind the contested findings f-001, f-002, f-004 and f-009. Resolve the m\_th fit inconsistency so the low-mass edge of the resonant window is set consistently, absorbing the structure-formation edge goal of d-04.; Quantify the upper (X-ray) and lower (asymmetry) edges of the viable band at 15, 20, 30, 40 and 60 keV, and assemble one viable-region table per asymmetry scenario.; Robustness check of the BBN conclusions: test whether 'BBN does not bind keV sterile-neutrino DM' survives other error budgets, and estimate the Yp shift from lepton asymmetry that the tool omits..

\textbf{Finding:} The tool grid (DW, m\_s = 15, 20, 30, 40, 60 keV; sin\textasciicircum{}2 2theta = 1e-15 to 1e-11; 45 points) gives no BBN exclusion anywhere: \textbar{}Delta chi\textasciicircum{}2\_BBN\textbar{} \textless{}= 0.0044 and Delta N\_eff \textless{}= 3.4e-4. Only 60 keV at 1e-11 overcloses under DW (Omega h\textasciicircum{}2 = 0.22). Lifetimes run from 9.68e27 s (15 keV, 1e-12) to 9.45e23 s (60 keV, 1e-11). The grid CSV has no X-ray exclusion column, so I took the X-ray flags from predict\_sterile\_nu\_xray\_line.

\textbf{Finding:} Tool X-ray flags: (15 keV, 1e-12) is excluded by NuSTAR M31 (Roach+2020) and XMM blank sky (Dessert+2020), with Gamma = 1.03e-28 s\textasciicircum{}-1. (30 keV, 1e-13; Gamma = 3.31e-28), (40 keV, 1e-14; Gamma = 1.39e-28) and (60 keV, 1e-14; Gamma = 1.06e-27 s\textasciicircum{}-1) are all reported 'not excluded by tabulated limits (within 30\% in mass)'. The tool's tabulated limits therefore do not reach the INTEGRAL/SPI or NuSTAR-GC regime above roughly 20-25 keV, so it cannot set the upper X-ray edge at 30-60 keV.

\textbf{Finding:} In the sterile-nu tool, the X-ray exclusion boundary at m\_s = 20 keV lies between sin\textasciicircum{}2 2theta = 1e-13 (not excluded, Gamma = 4.36e-29 s\textasciicircum{}-1) and 1e-12 (excluded by NuSTAR M31, d-02). At 30 keV the tool flags nothing even at sin\textasciicircum{}2 2theta = 1e-11 (Gamma = 3.31e-26 s\textasciicircum{}-1, E = 15 keV). The tool therefore has no tabulated X-ray limit for line energies above about 10-13 keV (m\_s \textgreater{}\textasciitilde{} 25 keV), which explains why it cannot reproduce the NuSTAR-GC / INTEGRAL upper edge at 30-60 keV.

\textbf{Finding:} The m\_th fit in 2507.18752 Eq. 8 has verbatim coefficients A = 1.329, a = 1.044, b = 0.055 (keV units), so at the paper's own benchmark (m\_s = 37 keV, sin\textasciicircum{}2 2theta = 1.5e-14) m\_th = 1.329 x 37\textasciicircum{}1.044 x (1.5e-14)\textasciicircum{}0.055 = 10.0 keV, not 8.6 keV. The f-002 value of 8.6 keV is an arithmetic error. The L fit (Eq. 9: C = 1.865e-12, alpha = -1.813, beta = -1.042) gives L = 0.68 at the same point, which confirms that part of f-002.

\textbf{Finding:} The consistent low-mass edge from m\_th \textgreater{} 8.5 keV (Zelko+, 95\% CL) uses the verified 2507.18752 Eq. 8: m\_s,min = [8.5/(1.329 (sin\textasciicircum{}2 2theta)\textasciicircum{}0.055)]\textasciicircum{}(1/1.044) keV. By hand this gives m\_s,min of about 36.5 keV at sin\textasciicircum{}2 2theta = 1e-15, 32.3 keV at 1e-14, 31.6 keV at 1.5e-14 and 28.6 keV at 1e-13. The L = 0.5 contour (Eq. 9) crosses the m\_th = 8.5 keV line near m\_s of about 31 keV and sin\textasciicircum{}2 2theta of about 2.8e-14. At sin\textasciicircum{}2 2theta = 1e-14 the edge point needs L of about 1.35. The paper's rounded summary 'm\_s \textgreater{}\textasciitilde{} 35 keV, sin\textasciicircum{}2 2theta \textless{}\textasciitilde{} 1e-14, L \textgreater{}\textasciitilde{} 0.5' (verbatim, Sec. IV) is consistent with this to within the fit and figure-reading precision; the X-ray exclusion, which I did not check, presumably cuts the larger-mixing corner. The board brief's 'm\_s \textgreater{}\textasciitilde{} 10 keV' is wrong and should read m\_s \textgreater{}\textasciitilde{} 30-37 keV (paper: \textgreater{}\textasciitilde{} 35 keV).

\textbf{Finding:} I verified the f-001/f-002 items in 2507.18752 verbatim. The fit validity range is '15 keV$\leq$m s$\leq$60 keV and 10\textasciicircum{}-15$\leq$sin2(2$\theta$)$\leq$3$\times$10\textasciicircum{}-13'; average errors are 1.3\% (m\_th) and 8.3\% (L). Further quotes: 'excluded at much greater than 95\% CL forL\textless{}0.1'; 'mth = 4.1+33.3 -1.5 keV [20]'; 'mth \textgreater{}8.4 keV (95\% CL) [18] andm th \textgreater{}8.5 keV (95\% CL) [19]'; Eq. 7 'mth,new$\approx$0.83m th,old + 0.36'; Fig. 7 caption: initial L 'given atT= 10 GeV' and 'reduced by a factor of 5 to 10 byT= 10 MeV'. Fig. 7 has m\_s axis ticks 15-60 keV and sin\textasciicircum{}2 2theta tick labels 1e-15, 1e-14 and 1e-13. This does not contradict the fit-validity upper end of 3e-13: the figure's labelled ticks and the fit's validation range are different quantities.

\textbf{Blocked --- Quantify the upper (X-ray) and lower (asymmetry) edges of the viable band at 15, 20, 30, 40 and 60 keV, and assemble one viable-region table per asymmetry scenario..} (The published limit curves (1609.00667 Figs. 7-8, 2507.20659 Fig. 1) cannot be extracted from HTML, and the tool has no X-ray limits above m\_s of about 25 keV, so I could not quantify the upper and lower edges at 30-60 keV.) I could not build the viable-region table, because the published X-ray edges and the zero-asymmetry lower edges are only available as plots I could not extract. What the tools did give is below.

\begin{enumerate}
  \item Tool grid (DW production; m\_s = 15, 20, 30, 40, 60 keV; sin\textasciicircum{}2 2theta from 1e-15 to 1e-11; 45 points; file /Users/nramachandra/hep-genesis-testing/projects/ramp-combo-sterile-bbn-xray/runs/261021\_143100/sterile\_nu\_grid\_DW\_677e35.csv):
\end{enumerate}
\begin{itemize}
  \item BBN excludes nothing. \textbar{}Delta chi\textasciicircum{}2\textbar{} is at most 0.0044 and Delta N\_eff at most 3.4e-4 (f-015).
  \item Only the point at 60 keV, sin\textasciicircum{}2 2theta = 1e-11 overcloses under DW (Omega h\textasciicircum{}2 = 0.22).
  \item The grid output has no X-ray flag column, so I ran the separate X-ray tool at selected points.
\end{itemize}
\begin{enumerate}
  \item Tool X-ray flags:
\end{enumerate}
\begin{itemize}
  \item 15 keV, sin\textasciicircum{}2 2theta = 1e-12: excluded by NuSTAR M31 and XMM blank sky (Gamma = 1.03e-28 s\textasciicircum{}-1).
  \item 20 keV: the tool's exclusion boundary lies between 1e-13 (not excluded, Gamma = 4.36e-29 s\textasciicircum{}-1) and 1e-12 (excluded, from d-02).
  \item 30 keV at 1e-13 and 1e-11, 40 keV at 1e-14, and 60 keV at 1e-14: none are flagged. At 30 keV, 1e-11 the rate is Gamma = 3.31e-26 s\textasciicircum{}-1.
  \item So the tool has no tabulated X-ray limits for line energies above about 10-13 keV (m\_s above about 25 keV) (f-016, f-017).
  \item The mismatch with the published limits comes from these missing limits, not from the decay-rate normalization, which matches 1.38e-32 (sin\textasciicircum{}2 2theta/1e-10)(m/keV)\textasciicircum{}5.
\end{itemize}
\begin{enumerate}
  \item Figure extraction failed for both papers (dead end x-02). 1609.00667 exposes only Figs. 3-4 in its HTML version, and 2507.20659 has no extractable figures. I have no digitized NuSTAR or INTEGRAL edge and no zero-asymmetry lower edge.
\end{enumerate}
What this leaves:

\begin{itemize}
  \item The success criteria are not met: I have no viable interval at five or more masses and no agreement within a factor of 2 with digitized curves, because no curves were digitized.
  \item Partial edges from the board: at 15 keV the upper edge is below 1e-12 (tool). The scenario with L around 0.5-1 is viable only at m\_s of about 35 keV or more and sin\textasciicircum{}2 2theta of about 1e-14 or less (f-001). The flavor-universal case with L around 1e-3 is excluded across 15-60 keV for full-DM production, since 2507.18752 excludes the window at L \textless{} 0.1.
  \item I did not address the Dessert vs Boyarsky blank-sky controversy beyond noting that the tool uses Dessert+2020.
\end{itemize}
Open: question q-06 asks for a PDF-based read-off of 1609.00667 Figs. 7-8 and 2507.20659 Fig. 1.

\textbf{Blocked --- Re-extract verbatim the numbers behind the contested findings f-001, f-002, f-004 and f-009. Resolve the m\_th fit inconsistency so the low-mass edge of the resonant window is set consistently, absorbing the structure-formation edge goal of d-04..} (The tool budget ran out before the 2510.24560, 2405.06509 and 1609.00667 text reads and a corrected cite() call could run.) The m\_th discrepancy is resolved: 10.0 keV is correct and the 8.6 keV in f-002 was an arithmetic error. The other contested papers were not re-checked because the budget ran out. The f-001 and f-002 numbers in 2507.18752 are now confirmed verbatim (tc-79, tc-75). The m\_th fit is Eq. 8: m\_th = 1.329 m\_s\textasciicircum{}1.044 (sin\textasciicircum{}2 2theta)\textasciicircum{}0.055, with masses in keV. The L fit is Eq. 9: L = 1.865e-12 m\_s\textasciicircum{}-1.813 (sin\textasciicircum{}2 2theta)\textasciicircum{}-1.042. Both are valid for 15-60 keV and sin\textasciicircum{}2 2theta from 1e-15 to 3e-13, with average errors of 1.3\% and 8.3\%. Also verbatim: m\_th \textgreater{} 8.5 keV (Zelko+, 95\% CL) and \textgreater{} 8.4 keV (galaxy counts); the Lyman-alpha value m\_th = 4.1 (+33.3/-1.5) keV; the conversion m\_th,new = 0.83 m\_th,old + 0.36 keV; the window being excluded at much more than 95\% CL for L \textless{} 0.1; and the conclusion 'opens up for L \textgreater{}\textasciitilde{} 0.5, m\_s \textgreater{}\textasciitilde{} 35 keV, sin\textasciicircum{}2 2theta \textless{}\textasciitilde{} 1e-14'. The Fig. 7 tick labels run from 1e-15 to 1e-13, while 3e-13 is the range over which the fits were validated. These are different quantities, so there is no conflict. At the 37 keV, sin\textasciicircum{}2 2theta = 1.5e-14 benchmark, my hand evaluation gives m\_th = 10.0 keV and L = 0.68 (f-018). Inverting the fit for m\_th = 8.5 keV gives the low-mass edge m\_s,min = [8.5/(1.329 (sin\textasciicircum{}2 2theta)\textasciicircum{}0.055)]\textasciicircum{}(1/1.044) keV. That is about 36.5 keV at 1e-15, 32.3 keV at 1e-14 and 28.6 keV at 1e-13 (my hand arithmetic, not a tool result). The L = 0.5 contour meets this edge near 31 keV and sin\textasciicircum{}2 2theta of about 2.8e-14 (f-019). The paper's rounded \textgreater{}\textasciitilde{} 35 keV is consistent with this. The brief's 'm\_s \textgreater{}\textasciitilde{} 10 keV' is wrong and should read m\_s \textgreater{}\textasciitilde{} 30-37 keV, or \textgreater{}\textasciitilde{} 35 keV as the paper states (f-019, f-020). Still open (dead end x-03): verbatim re-extraction for f-004 (2405.06509), f-009 (XRISM 2510.24560; the 0-30000 chunk held only the author list) and f-011 (1609.00667). These are 'not found', not refuted. The cite() of the 2507.18752 quotes failed on one line-broken quote, so references.bib has no new entries from this direction. The tc-42 misattribution was not corrected.

\textbf{Finding:} plot\_sterile\_nu\_bbn\_impact (Figure 12, bbn\_vs\_neff.png) gives an SM baseline of Yp = 0.24682 and D/H = 2.4433e-5, with slopes dYp/dN\_eff = 0.01351 and d(D/H x1e5)/dN\_eff = 0.3308 near the SM. One fully thermalised sterile neutrino (N\_eff = 4.044) gives Yp = 0.2596 and D/H = 2.777e-5. The tool's BBN-only limit is Delta N\_eff \textless{} 0.548 (95\%). The plot does not depend on m\_s: it is a function of N\_eff only, so there is no separate 7.1 keV run.

\textbf{Finding:} Correction to f-007 (my own reconstruction from the tool outputs): the tool's BBN chi\textasciicircum{}2 uses Yp and D/H only, not Planck N\_eff. Its D/H error is about 0.065e-5 in total, which implies a D/H theory error of about 0.060e-5 added in quadrature to the observational 0.025e-5. With that error the SM D/H offset is -1.6 sigma, not -4.1 sigma; -4.1 sigma holds only with the observational error alone. Check: baseline chi\textasciicircum{}2 = 0.23 (Yp) + 2.54 (D/H) = 2.77. This reproduces the tool's Delta chi\textasciicircum{}2 = -1.14 at Delta N\_eff = 0.161, Delta chi\textasciicircum{}2 = 5.9 at Delta N\_eff = 0.62, and the Delta N\_eff \textless{} 0.548 limit.

\textbf{Finding:} The claim that BBN does not bind keV DW sterile-neutrino DM holds under every error budget I tested (own linearized estimates, using the tool slopes and the linear DW scaling Delta N\_eff = 6.18e6 sin\textasciicircum{}2 2theta at 7.1 keV and 2.95e7 sin\textasciicircum{}2 2theta at 50 keV taken from the grid). The 95\% limits are:

\begin{itemize}
  \item Yp only: Delta N\_eff \textless{} 0.455
  \item Tool default (Yp plus D/H with theory error): \textless{} 0.548
  \item Yp plus D/H with observational error only (SM chi\textasciicircum{}2 = 17.4): \textless{} 0.61
  \item D/H only with theory error: \textless{} 0.81
  \item Planck N\_eff 2.99 +/- 0.17, as an extreme check (a keV sterile neutrino is non-relativistic at recombination): \textless{} 0.28
These translate to sin\textasciicircum{}2 2theta \textless{} 4.5e-8 to 1.3e-7 at 7.1 keV and 9.5e-9 to 2.8e-8 at 50 keV, which is 180-510 times (7.1 keV) and 1250-3600 times (50 keV) above the DW relic line (2.545e-10 and 7.583e-12). Every budget therefore clears the 10x criterion. The tool grid agrees: at 50 keV, sin\textasciicircum{}2 2theta = 2.6e-8 gives Delta N\_eff = 0.766 (excluded), while 1e-8 is allowed. At 7.1 keV, 7.5e-8 is allowed (Delta chi\textasciicircum{}2 = 1.83).

\end{itemize}
\textbf{Finding (preliminary):} Own estimate, not a tool result: I used Delta Yp \textasciitilde{} -0.23 xi\_e (literature-standard sensitivity) with xi \textasciitilde{} 1.46 L (asymmetry normalized to photons, T\_nu = T\_gamma before e+e- annihilation; this conversion depends on convention at the factor-of-4 level).

\begin{itemize}
  \item Flavor-universal L = 1e-3 gives xi\_e \textasciitilde{} 1.5e-3 and Delta Yp \textasciitilde{} -3.4e-4, or 0.08 sigma against sigma\_Yp = 0.004. Negligible.
  \item A 1 sigma (2 sigma) Yp shift needs a residual xi\_e \textasciitilde{} 0.017 (0.035), i.e. a residual L\_e \textasciitilde{} 0.012 (0.024).
  \item 2507.20659 (L\_e = -L\_mu = 0.1): without equilibration, xi\_e \textasciitilde{} 0.15 and Delta Yp \textasciitilde{} -0.034 (about 8 sigma). Their window therefore needs the residual xi\_e at BBN to be below about 0.03, which zero-total-asymmetry equilibration is meant to deliver.
  \item 2507.18752 (L \textasciitilde{} 0.68 at 20 MeV, down 5-10 times by 10 MeV): if the residual asymmetry sat in nu\_e, xi\_e \textasciitilde{} 0.1-0.2 would give Delta Yp \textasciitilde{} -0.023 to -0.046 (6-12 sigma). BBN viability therefore rests on the asymmetry staying out of nu\_e (mostly L\_mu), and the tool cannot test this.
  \item The xi-induced Delta N\_eff (\textasciitilde{}0.22 xi\textasciicircum{}2 per flavor, \textasciitilde{}0.002 at xi = 0.1) is negligible next to the Yp effect.
\end{itemize}
\textbf{Completed --- Robustness check of the BBN conclusions: test whether 'BBN does not bind keV sterile-neutrino DM' survives other error budgets, and estimate the Yp shift from lepton asymmetry that the tool omits..} 1. BBN impact plot (Figure 12, /Users/nramachandra/hep-genesis-testing/projects/ramp-combo-sterile-bbn-xray/runs/261021\_143100/bbn\_vs\_neff.png; tc-...-86). The plot depends only on N\_eff, not on m\_s, so no separate 7.1 keV run was needed. SM baseline: Yp = 0.24682, D/H = 2.4433e-5. Slopes near the SM: dYp/dN\_eff = 0.01351 and d(D/H x1e5)/dN\_eff = 0.3308. One thermalised sterile neutrino (N\_eff = 4.044) gives Yp = 0.2596 and D/H = 2.777e-5. The tool's BBN-only limit is Delta N\_eff \textless{} 0.548 (95\%) (f-021).

\begin{enumerate}
  \item Implied D/H theory error and the f-007 correction (f-022). The tool's chi\textasciicircum{}2 uses Yp and D/H only. My reconstruction gives a total D/H error of about 0.065e-5, which implies a theory error of about 0.060e-5 added in quadrature to the observational 0.025e-5. The baseline then splits as chi\textasciicircum{}2 = 0.23 (Yp) + 2.54 (D/H) = 2.77, and the reconstruction reproduces the tool's Delta chi\textasciicircum{}2 values at Delta N\_eff = 0.161 and 0.62. Correction to f-007: the SM D/H offset is -1.6 sigma against the tool's own error budget. The -4.1 sigma figure holds only with the observational error alone.
\end{enumerate}
\begin{enumerate}
  \item Robustness under other error budgets (f-023; own linearized estimates, partly checked against the new tool grid sterile\_nu\_grid\_DW\_7759e2.csv, tc-...-87). 95\% limits on Delta N\_eff:
\end{enumerate}
\begin{itemize}
  \item Yp only: \textless{} 0.455
  \item Tool default: \textless{} 0.548
  \item Yp plus D/H with observational error only: \textless{} 0.61
  \item D/H with theory error only: \textless{} 0.81
  \item Planck N\_eff (extreme check): \textless{} 0.28
In mixing angle these are sin\textasciicircum{}2 2theta \textless{} 4.5e-8 to 1.3e-7 at 7.1 keV and 9.5e-9 to 2.8e-8 at 50 keV. That is 180-510 times (7.1 keV) and 1250-3600 times (50 keV) above the DW relic line. The more-than-10x criterion is met under every budget, so the claim that BBN does not bind keV DW sterile-neutrino DM holds.

\end{itemize}
\begin{enumerate}
  \item Yp shift from lepton asymmetry (f-024; own estimate using Delta Yp \textasciitilde{} -0.23 xi\_e and xi \textasciitilde{} 1.46 L, convention-dependent, low confidence).
\end{enumerate}
\begin{itemize}
  \item Flavor-universal L = 1e-3: Delta Yp \textasciitilde{} -3.4e-4, or 0.08 sigma against sigma\_Yp = 0.004. Negligible.
  \item A 2 sigma shift needs a residual xi\_e of about 0.035.
  \item Unequilibrated benchmarks would give Delta Yp of about -0.034 for 2507.20659 (L\_e = 0.1) and -0.02 to -0.05 for 2507.18752 if the asymmetry sat in nu\_e. Both are about 6-12 sigma, so those windows depend on the residual xi\_e at BBN being below about 0.03.
\end{itemize}
Open question q-07: what residual xi\_e do the benchmarks actually reach after flavor equilibration?

\textbf{Critic, wave 2} --- continue: 4 findings stand up to the evidence, 6 questioned ('Tool X-ray flags: (15 keV, 1e-12) is excluded by NuSTAR M31 (Roach+2020) and XMM blank sky': The quoted excerpts of tc-80 to tc-83 match the decay rates and lifetimes, and they agree with Gamma = 1.38e-32 (sin\textasciicircum{}2 2theta/1e-10)(m/keV)\textasciicircum{}5: about 1.41e-28 at 40 keV and 1.07e-27 at 60 keV. Every excerpt is cut off at 'stacked\_', before the exclusion verdict, so the claimed 'excluded by NuSTAR M31 / XMM blank sky' and 'not excluded by tabulated limits (within 30\% in mass)' strings do not appear in the evidence. The finding also puts the tool's coverage limit at about 20-25 keV, while f-017 puts it at line energies of about 10-13 keV (m\_s about 25 keV). The two thresholds need to be reconciled.; 'In the sterile-nu tool, the X-ray exclusion boundary at m\_s = 20 keV lies between sin\textasciicircum{}2 2t': The same truncation applies. tc-91 and tc-92 show Gamma = 4.36e-29 and 3.31e-26 s\textasciicircum{}-1, and both values are correct. The 'not excluded' verdicts at 20 keV/1e-13 and 30 keV/1e-11 are not in the quoted text. The 'excluded at 1e-12' half of the 20 keV bracket rests on a d-02 tool call that is not cited here. A coverage gap in the tool is plausible, but the evidence as quoted does not establish it.; 'The m\_th fit in 2507.18752 Eq. 8 has verbatim coefficients A = 1.329, a = 1.044, b = 0.055': The arithmetic checks out: 1.329 x 37\textasciicircum{}1.044 x (1.5e-14)\textasciicircum{}0.055 = 10.0 keV, and the L fit gives 0.68. The only evidence cited is tc-79, and its excerpt is text about Fig. 4 and the asymmetry discussion. It does not contain Eq. 8, Eq. 9 or the coefficients 1.329/1.044/0.055 that the finding calls verbatim. This is the third wave in which the 2507.18752 fit coefficients appear without a matching quoted excerpt.; 'The consistent low-mass edge from m\_th \textgreater{} 8.5 keV (Zelko+, 95\% CL) uses the verified 2507.1': Hand arithmetic reproduces the inversion: 36.5, 32.3 and 28.6 keV at 1e-15, 1e-14 and 1e-13, an L = 0.5 crossing near 30.6 keV and 2.8e-14, and L of about 1.33 at the 1e-14 edge. The finding still inherits f-018's unquoted coefficients. It also leaves two points unchecked. First, it does not say whether the m\_th \textgreater{} 8.5 keV bound (Zelko+) uses the same thermal-mass convention as Eq. 8; if it is an 'old' value, Eq. 7 maps it to about 7.4 keV and the edge moves down by about 4-5 keV. Second, the crossing at sin\textasciicircum{}2 2theta of about 2.8e-14 lies above the paper's quoted sin\textasciicircum{}2 2theta \textless{}\textasciitilde{} 1e-14, which implies an X-ray cut near 31 keV that the finding has not verified.).

Issues: The board brief still states 'viable regions at m\_s \textgreater{}\textasciitilde{} 10 keV' and 'no such tools were available to me'. Both are now wrong: the low-mass edge is about 30-37 keV (paper: \textgreater{}\textasciitilde{} 35 keV), and the sterile-nu tools were used extensively. The final synthesis must correct the brief.; The m\_th convention is unchecked. Eq. 7 (m\_th,new = 0.83 m\_th,old + 0.36) shows the paper distinguishes two conventions. The low-mass edge moves by several keV depending on which convention the 8.5 keV bound and the Eq. 8 output use.; There is still no quantitative X-ray upper edge at m\_s = 30-60 keV. This is exactly the mass range where the only viable full-DM window (f-001) sits, so the central question of which region is viable is answered only by quoting 2507.18752's rounded summary, not by an independent cross-check..

Stopping after wave 2: budget reached (2/2 waves used). Writing the report from the board.

\textit{Research report:}\par
\subsubsection*{Sterile-neutrino dark matter in 2024-2026: combining BBN and X-ray constraints on the (m\_s, sin\textasciicircum{}2 2theta) plane}
\subsubsection*{Abstract}
This investigation combined sterile-neutrino tool calculations with a reading of recent arXiv papers to ask which (m\_s, sin\textasciicircum{}2 2theta) regions remain viable for keV sterile-neutrino dark matter.

\textbf{BBN.} The tool's BBN module responds only to Delta N\_eff. On the non-resonant (Dodelson-Widrow, DW) relic-density line it gives Delta N\_eff between 2.2e-4 and 1.6e-3. The resulting abundance shifts are negligible (\textbar{}Delta chi\textasciicircum{}2\textbar{} of 0.02 or less). BBN begins to exclude only at mixing angles 180-3600 times above the relic line, under every error budget tested. BBN therefore does not constrain keV sterile-neutrino dark matter through Delta N\_eff. Where BBN does matter is through the lepton asymmetry that resonant production requires, and the tool does not model that effect.

\textbf{X-rays.} The tool's decay rates agree with the standard radiative formula to within about 2\%. Its tabulated limits exclude DW production as all of the dark matter for m\_s = 7-20 keV. The newest XRISM cluster-stack limit is sin\textasciicircum{}2 2theta below about 4-5e-10 at 7.1 keV. That does not exclude the 3.5 keV line, which implies about 7e-11.

\textbf{What survives.} For sterile neutrinos to be all of the dark matter, the literature (Vogel et al., arXiv:2507.18752) gives a resonant (Shi-Fuller) window. It requires an initial lepton asymmetry L of about 0.5 or more, m\_s of about 35 keV or more, and sin\textasciicircum{}2 2theta of about 1e-14 or less. Two other cases:

\begin{itemize}
  \item A flavor-universal asymmetry of about 1e-3 is excluded for full dark matter across 15-60 keV.
  \item A zero-total-asymmetry flavor configuration (arXiv:2507.20659) widens the viable mixing range, but its edges could not be quantified here.
\end{itemize}
The quantitative upper (X-ray) and lower (asymmetry) edges at 30-60 keV remain open. They exist only as plots that could not be extracted.

\textbf{Correction to the brief.} The initial survey brief stated a low-mass edge of ''m\_s \textgreater{}\textasciitilde{} 10 keV'' and said that no BBN tools were available. Both statements are wrong. The edge is about 30-37 keV, and the sterile-neutrino tools were used throughout.

\subsubsection*{Background}
\textbf{X-ray searches.} The XRISM Collaboration stacked 3.75 Ms of Resolve data on ten galaxy clusters (arXiv:2510.24560) and found no unidentified lines. For a 7.1 keV particle, its 3-sigma decay-rate limit is Gamma \textless{} 0.97e-27 s\textasciicircum{}-1. The paper also gives a mixing-angle fit, valid for line energies of 2.5-15 keV:

sin\textasciicircum{}2 2theta / 1e-10 \textless{} 4.1 (E/3.55)\textasciicircum{}-4 f1(E)

The NuSTAR Galactic Center analysis (arXiv:1609.00667) improved earlier limits by about an order of magnitude near m\_s = 20 keV and implied m\_s \textless{}\textasciitilde{} 16 keV in the nuMSM. It quotes the decay-rate normalization used throughout this report:

Gamma = 1.38e-32 s\textasciicircum{}-1 (sin\textasciicircum{}2 2theta/1e-10)(m/keV)\textasciicircum{}5

\textbf{Lepton asymmetries and BBN.} Froustey and Pitrou (arXiv:2405.06509) used full three-flavor neutrino transport with BBN. They find that the space of allowed initial lepton asymmetries is ''generically unbound'' with current abundance and CMB data, because flavors do not fully equilibrate. A tight bound of about 1e-3 holds only for a flavor-universal asymmetry. A quantum-kinetic study of flavor equilibration (arXiv:2502.14960) supports the same picture: large flavor asymmetries can evade BBN and CMB limits.

\textbf{Resonant-production windows.}

\begin{itemize}
  \item Vogel et al. (arXiv:2507.18752) revive Shi-Fuller production with L of order 1 above about 20 MeV. They combine NuSTAR and INTEGRAL/SPI X-ray limits with structure-formation bounds and provide fitting functions for the thermal-equivalent mass m\_th and for the required L.
  \item Akita, Hamaguchi and Ovchynnikov (arXiv:2507.20659) assume large flavor asymmetries with zero total asymmetry, washed out by oscillations at about 15 MeV. They find that the viable mixing range grows by up to two orders of magnitude for m\_s up to 60 keV, and they release a public code (sterile-dm-lfa).
  \item Gorbunov et al. (arXiv:2502.17374) argue that unequal flavor asymmetries soften BBN limits and give colder spectra.
  \item Two further papers were listed in the survey but not checked. arXiv:2402.11902 (an Affleck-Dine source for the asymmetry) was not read. arXiv:2609.31340 (saturation of resonant production at high asymmetry) was never checked against arXiv metadata, and its identifier looks implausible. Neither is relied on below.
\end{itemize}
\subsubsection*{Results}
\paragraph{What the sterile-neutrino tools can and cannot compute}
\textbf{BBN module.} BBN is a PRyMordial-calibrated fit for Yp, D/H and Li7. It is driven only by Delta N\_eff and is valid for m\_s below about 100 keV. A point is excluded at Delta chi\textasciicircum{}2 \textgreater{} 3.84 against two measurements:

\begin{itemize}
  \item Yp = 0.2449 +/- 0.004 (Aver+ 2021)
  \item D/H = 2.547 +/- 0.025 e-5 (Cooke+ 2018)
\end{itemize}
Earlier the criterion was described as also including Planck N\_eff = 2.99 +/- 0.17. A later reconstruction shows the chi\textasciicircum{}2 actually uses Yp and D/H only (see the error-budget subsection below).

\textbf{Production mechanisms.} DW production is the only mechanism the tool models reliably. Its Shi-Fuller mode accepts a lepton\_asymmetry argument, but:

\begin{itemize}
  \item Its Delta N\_eff is a labelled heuristic (0.1 times the DW value).
  \item At m\_s = 7.1 keV, sin\textasciicircum{}2 2theta = 7e-11 and L = 1e-3 it returned Omega h\textasciicircum{}2 = 432, which is unphysical.
  \item The asymmetry's own effect on BBN (the shift in n/p from xi\_e) is not modelled.
\end{itemize}
\textbf{Shi-Fuller benchmark check.} At the 2507.18752 benchmark (m\_s = 37 keV, sin\textasciicircum{}2 2theta = 1.5e-14, L = 0.1, the tool's maximum) the tool gives:

\begin{itemize}
  \item Delta N\_eff = 3.5e-8
  \item Yp = 0.24682 and D/H = 2.443e-5, identical to the Standard Model values
  \item BBN Delta chi\textasciicircum{}2 = 0
  \item Decay rate 1.42e-28 s\textasciicircum{}-1 at E = 18.5 keV
  \item Omega\_s h\textasciicircum{}2 = 25.1
\end{itemize}
That relic density is about 200 times the observed dark matter. The paper's fit instead places this point at the observed density for L of about 0.7. The tool therefore cannot test the published resonant windows, either on the BBN side or on the relic-density side.

\paragraph{DW scan of the (m\_s, sin\textasciicircum{}2 2theta) plane}
The scan covered m\_s = 1-100 keV and sin\textasciicircum{}2 2theta = 1e-14 to 1e-6 on a 60x60 grid (files scan\_DW\_9bfc3d.npz/.csv). Of the 3600 points, 546 are BBN-excluded and 1740 overclose.

The BBN 95\% upper limits on sin\textasciicircum{}2 2theta, by mass, are:

\begin{longtable}{@{}p{0.137\textwidth}p{0.137\textwidth}p{0.137\textwidth}p{0.137\textwidth}p{0.137\textwidth}p{0.137\textwidth}p{0.137\textwidth}@{}}
\toprule
\textbf{m\_s (keV)} & \textbf{1} & \textbf{7} & \textbf{13} & \textbf{26} & \textbf{50} & \textbf{100} \\ \midrule \endhead
BBN limit on sin\textasciicircum{}2 2theta & 3.9e-7 & 8.2e-8 & 4.4e-8 & 2.4e-8 & 1.7e-8 & 9.3e-9 \\
\bottomrule\end{longtable}
These limits are two or more orders of magnitude weaker than the overclosure boundary.

[[FIGURE 10]]

Figure 10 shows the DW exclusion plane. The Omega h\textasciicircum{}2 = 0.12 line runs from sin\textasciicircum{}2 2theta = 2.5e-10 at 7.1 keV down to 7.6e-12 at 50 keV, with overclosure above it. The BBN-excluded region lies only above about 1e-8 to 1e-7 and so never reaches the relic line.

The tool also tabulated BBN numbers at 84 points. The precise DW relic-line points, with their BBN output, are:

\begin{longtable}{@{}p{0.240\textwidth}p{0.240\textwidth}p{0.240\textwidth}p{0.240\textwidth}@{}}
\toprule
\textbf{m\_s (keV)} & \textbf{sin\textasciicircum{}2 2theta} & \textbf{Delta N\_eff} & \textbf{Delta chi\textasciicircum{}2\_BBN} \\ \midrule \endhead
7.1 & 2.545e-10 & 1.6e-3 & -0.020 \\
10 & 1.374e-10 & 1.1e-3 & -0.014 \\
20 & 3.946e-11 & 5.6e-4 & -0.007 \\
50 & 7.583e-12 & 2.2e-4 & -0.003 \\
\bottomrule\end{longtable}
At 7.1 keV and sin\textasciicircum{}2 2theta = 3e-10 the tool gives Yp = 0.246840 and D/H = 2.44395e-5. Both are indistinguishable from the Standard Model baseline (Yp = 0.246815, D/H = 2.4433e-5, N\_eff = 3.044).

BBN starts to exclude only far above the relic line:

\begin{itemize}
  \item At 7.1 keV, near sin\textasciicircum{}2 2theta of about 1e-7: Delta N\_eff = 0.62, Yp = 0.2549, D/H = 2.649e-5, Delta chi\textasciicircum{}2 = 5.9.
  \item At 50 keV, near 3e-8: Delta N\_eff = 0.88, Delta chi\textasciicircum{}2 = 16.8.
\end{itemize}
A second grid at m\_s = 15, 20, 30, 40 and 60 keV, with sin\textasciicircum{}2 2theta from 1e-15 to 1e-11 (45 points), found no BBN exclusion anywhere: \textbar{}Delta chi\textasciicircum{}2\textbar{} was at most 0.0044 and Delta N\_eff at most 3.4e-4. Only the point at 60 keV with sin\textasciicircum{}2 2theta = 1e-11 overcloses under DW (Omega h\textasciicircum{}2 = 0.22).

\paragraph{BBN response to N\_eff and robustness of the error budget}
[[FIGURE 12]]

Figure 12 shows the tool's BBN predictions as a function of N\_eff only; m\_s does not enter. Its key numbers:

\begin{itemize}
  \item Standard Model baseline: Yp = 0.24682, D/H = 2.4433e-5.
  \item Slopes near the Standard Model: dYp/dN\_eff = 0.01351 and d(D/H x 1e5)/dN\_eff = 0.3308.
  \item One fully thermalised sterile neutrino (N\_eff = 4.044): Yp = 0.2596 and D/H = 2.777e-5.
  \item BBN-only limit: Delta N\_eff \textless{} 0.548 (95\%).
\end{itemize}
\textbf{Reconstructing the tool's error budget.} The tool's numbers are reproduced only if its chi\textasciicircum{}2 uses Yp and D/H, with a D/H theory error of about 0.060e-5 added in quadrature to the observational 0.025e-5. The baseline then splits as chi\textasciicircum{}2 = 0.23 (Yp) + 2.54 (D/H) = 2.77. This reconstruction matches the tool's Delta chi\textasciicircum{}2 at three points:

\begin{itemize}
  \item Delta chi\textasciicircum{}2 = -1.14 at Delta N\_eff = 0.161
  \item Delta chi\textasciicircum{}2 = 5.9 at Delta N\_eff = 0.62
  \item Delta chi\textasciicircum{}2 = -0.367 at Delta N\_eff = 0.031; adding Planck N\_eff to the chi\textasciicircum{}2 would give about -0.21 instead
\end{itemize}
It follows that the Standard Model D/H offset is -1.6 sigma against the tool's own error budget. The earlier figure of -4.1 sigma holds only if the observational error is used alone.

\textbf{Alternative budgets.} These are my own linearized estimates, using the grid scaling Delta N\_eff = 6.18e6 sin\textasciicircum{}2 2theta at 7.1 keV and 2.95e7 sin\textasciicircum{}2 2theta at 50 keV. The 95\% limits on Delta N\_eff are:

\begin{longtable}{@{}p{0.480\textwidth}p{0.480\textwidth}@{}}
\toprule
\textbf{Error budget} & \textbf{95\% limit on Delta N\_eff} \\ \midrule \endhead
Yp only & \textless{} 0.455 \\
Tool default (Yp plus D/H with theory error) & \textless{} 0.548 \\
Yp plus D/H, observational error only & \textless{} 0.61 \\
D/H only, with theory error & \textless{} 0.81 \\
Planck N\_eff (extreme check) & \textless{} 0.28 \\
\bottomrule\end{longtable}
The Planck case is an extreme check only, because a keV sterile neutrino is non-relativistic at recombination.

In mixing angle, these limits correspond to:

\begin{itemize}
  \item sin\textasciicircum{}2 2theta \textless{} 4.5e-8 to 1.3e-7 at 7.1 keV, which is 180-510 times above the relic line;
  \item sin\textasciicircum{}2 2theta \textless{} 9.5e-9 to 2.8e-8 at 50 keV, which is 1250-3600 times above the relic line.
\end{itemize}
A direct tool grid agrees. At 50 keV, sin\textasciicircum{}2 2theta = 2.6e-8 gives Delta N\_eff = 0.766 and is excluded, while 1e-8 is allowed. At 7.1 keV, 7.5e-8 is allowed (Delta chi\textasciicircum{}2 = 1.83).

\paragraph{Lepton-asymmetry shift of Yp (estimate only)}
The tool omits the shift in Yp caused by the asymmetry itself. I estimated it by hand, using Delta Yp of about -0.23 xi\_e and xi of about 1.46 L. These are low-confidence numbers, because the L-to-xi conversion is convention-dependent at the factor-of-4 level.

\begin{itemize}
  \item \textbf{Flavor-universal L = 1e-3.} Delta Yp is about -3.4e-4, or 0.08 sigma. This is negligible.
  \item \textbf{Threshold for a visible shift.} A 2-sigma shift in Yp needs a residual xi\_e of about 0.035.
  \item \textbf{2507.20659 benchmark (L\_e = -L\_mu = 0.1).} Without equilibration, Delta Yp would be about -0.034 (about 8 sigma).
  \item \textbf{2507.18752 benchmark.} If its asymmetry sat in nu\_e, Delta Yp would be about -0.02 to -0.05 (6-12 sigma).
\end{itemize}
Both resonant windows are therefore BBN-safe only if the residual xi\_e at about 1 MeV stays below about 0.03. This is what flavor equilibration (for zero total asymmetry) or keeping the asymmetry mostly in L\_mu is meant to deliver. The tool cannot check it.

\paragraph{X-ray constraints}
\textbf{Normalization.} The tool's decay rates follow the standard formula with a 1.36e-29 prefactor:

\begin{itemize}
  \item On the DW relic line they are 6.25e-28 s\textasciicircum{}-1 (7.1 keV), 1.87e-27 s\textasciicircum{}-1 (10 keV), 1.72e-26 s\textasciicircum{}-1 (20 keV) and 3.23e-25 s\textasciicircum{}-1 (50 keV), within about 1.5\% of the analytic values.
  \item At 7.1 keV with sin\textasciicircum{}2 2theta = 4e-10 the tool gives 9.82e-28 s\textasciicircum{}-1, within 2\% of the XRISM limit of 0.97e-27 s\textasciicircum{}-1.
  \item At 20 keV with sin\textasciicircum{}2 2theta = 1e-12 it gives 4.36e-28 s\textasciicircum{}-1, against 4.42e-28 s\textasciicircum{}-1 from the NuSTAR Eq. 2.
\end{itemize}
The tool's fluxes are whole-halo values with no field-of-view cut, so only decay rates, not fluxes, can be compared with published limits.

\textbf{XRISM limits.} Evaluating the XRISM fit by hand gives the following upper limits on sin\textasciicircum{}2 2theta:

\begin{longtable}{@{}p{0.160\textwidth}p{0.160\textwidth}p{0.160\textwidth}p{0.160\textwidth}p{0.160\textwidth}p{0.160\textwidth}@{}}
\toprule
\textbf{m\_s (keV)} & \textbf{6} & \textbf{7.1} & \textbf{8} & \textbf{10} & \textbf{14} \\ \midrule \endhead
XRISM limit on sin\textasciicircum{}2 2theta & about 1.3e-9 & about 4-5e-10 & about 2.3e-10 & about 5.7e-11 & about 1.2e-11 \\
\bottomrule\end{longtable}
At 7.1 keV the two routes differ slightly: the decay-rate limit gives 3.9e-10 and the polynomial fit gives 4.8e-10. The 3.5 keV line implies a rate of about 2e-28 s\textasciicircum{}-1 (sin\textasciicircum{}2 2theta of about 7e-11), a factor of about 5 below the XRISM limit, so XRISM does not exclude the line. These XRISM numbers were not re-verified verbatim in a second pass; that attempt hit the author-list chunk of the paper. They should be treated as single-read values.

\textbf{The tool's tabulated limits} (NuSTAR M31, Roach+ 2020; XMM blank sky, Dessert+ 2020):

\begin{itemize}
  \item They exclude the DW relic points at 7.1 and 10 keV (both limits) and at 20 keV (NuSTAR M31). DW as all of the dark matter is therefore X-ray-excluded for 7-20 keV.
  \item At 7.1 keV they also exclude sin\textasciicircum{}2 2theta = 4e-10, so these limits are stronger than XRISM there.
  \item At 15 keV, sin\textasciicircum{}2 2theta = 1e-12 is excluded.
  \item At 20 keV the exclusion boundary lies between 1e-13 (allowed, Gamma = 4.36e-29 s\textasciicircum{}-1) and 1e-12.
\end{itemize}
Above m\_s of about 25 keV (line energies above about 10-13 keV) the tool has no tabulated limits. Nothing is flagged at 30 keV even at sin\textasciicircum{}2 2theta = 1e-11 (Gamma = 3.31e-26 s\textasciicircum{}-1), nor at 40 or 60 keV with sin\textasciicircum{}2 2theta = 1e-14. The 50 keV DW relic point is ''not excluded'' only because of this gap. Published NuSTAR Galactic Center, INTEGRAL and Lyman-alpha limits exclude DW dark matter at 50 keV. These high-mass tool results describe the tool's coverage, not the physics.

\textbf{NuSTAR Galactic Center.} 1609.00667 gives its sin\textasciicircum{}2 2theta limits only as curves (Figs. 7-8). Extracting those figures failed, so no numerical NuSTAR or INTEGRAL edge is available at 30-60 keV.

\paragraph{Resonant-production window from the literature}
The following items from 2507.18752 were checked against the full text in a second pass:

\begin{itemize}
  \item The window opens for L \textgreater{}\textasciitilde{} 0.5 (given at T = 10 GeV), m\_s \textgreater{}\textasciitilde{} 35 keV and sin\textasciicircum{}2 2theta \textless{}\textasciitilde{} 1e-14.
  \item It is excluded at much more than 95\% CL for L \textless{} 0.1.
  \item Structure
\end{itemize}
\end{document}
